\section{Related Work}
\label{sec:related}

\subsection{LLMs for FHIR and Clinical Records}

The most direct predecessor to this work is LLMonFHIR \citep{schmiedmayer2025llmonfhir},
a Stanford physician-validated iOS application that allows patients to query their
own FHIR records using natural language. The JACC publication evaluates 210 physician-rated
responses on Synthea patients and reports median accuracy scores of 5/5, but explicitly
identifies context size as the central unresolved technical challenge: raw FHIR bundles
can reach 800,000 lines of JSON for a single patient, far exceeding current LLM context
windows. The earlier arXiv version of the same project \citep{schmiedmayer2024llmonfhir}
addresses this with GPT-4 function-calling to select relevant resource types before
sending data to the model. Our work takes a different angle: given a bounded,
task-relevant subset of FHIR resources (only MedicationRequest), how does the
\textit{format} of that subset affect model performance? This is the layer below
resource selection that prior work has not systematically examined.

\citet{wu2026plannerauditor} build a dual-agent discharge planning system using MIMIC-IV
FHIR data, where a Planner LLM generates structured discharge action plans and a
deterministic Auditor validates coverage and calibration. Their self-improvement loop
raises task coverage from 32\% to 86\%. This work is relevant because it shows that
FHIR-native LLM pipelines can reach high task coverage when structured validation is
applied, a design pattern that could be combined with our serialisation findings
for production deployment.

\subsection{Medication Extraction and FHIR-Specific Tasks}

\citet{li2024fhirgpt} introduce FHIR-GPT, which converts free-text clinical notes
and discharge summaries into FHIR MedicationStatement resources using GPT-4 and
few-shot prompting. On 3,671 medication snippets from the n2c2 2018 dataset, GPT-4
achieves over 90\% exact match, substantially outperforming prior NLP pipelines.
This is the closest prior work to our task: LLMs applied to medication-specific
FHIR resources. The key distinction is directionality. FHIR-GPT extracts from
unstructured clinical text \textit{into} FHIR format; we extract from FHIR format
\textit{into} a medication list. The inverse task uses the same data format but
requires different model capabilities: structured input comprehension rather than
structured output generation.

\subsection{Hallucination, Omission, and Clinical Safety}

\citet{asgari2025clinicalsafety} conduct the largest manual evaluation of LLM clinical
note generation to date (12,999 annotated sentences) and measure hallucination at 1.47\%
and omission at 3.45\%. They find that iterative prompt engineering can bring error
rates below human baseline and identify omission as a persistent challenge even after
optimisation. This paper establishes the methodological framing that we extend:
treating precision failures (hallucination) and recall failures (omission) as distinct
error types with different clinical consequences.

\citet{zeba2025granularfactcheck} introduce a proposition-level, LLM-free fact-checking
module that validates clinical text against EHR source data using discrete logical checks
(negation, temporal consistency, numerical comparison). Their approach achieves F1 = 0.856
for hallucination detection and demonstrates that systematic validation is tractable
without another LLM in the loop. The distinction between hallucination and omission
they formalise maps directly onto our finding that omission dominates hallucination in
this benchmark: mean precision exceeds mean recall in every model--strategy combination,
even though precision failures do occur.

\citet{kim2025medicalhallucination} evaluate medical hallucination across 11 foundation
models on 7 clinical tasks, with survey input from 70 clinicians. Their finding that
general-purpose models outperform medical-specialised ones (76.6\% vs.\ 51.3\%
hallucination-free) and that hallucination is a reasoning failure rather than a knowledge
deficit aligns closely with our BioMistral result: domain pretraining without instruction
tuning produces a model that generates no useful output on a structured extraction task.

\subsection{Prompting and Format Effects in Clinical NLP}

\citet{delaunay2025fhirprompting} evaluate three LLMs on converting Spanish clinical
reports into FHIR bundles using six prompting strategies. They find that two-step
prompting (extract then format) significantly outperforms single-step prompting, and
that temperature 0 is optimal for FHIR compliance. This is the most directly comparable
prompting study to our work. We differ in task direction (FHIR comprehension vs.\
FHIR generation), language, model family, and the specific format variable we test,
but the shared finding is that how you structure the model's task matters as much as
which model you choose.

\citet{kruse2025temporal} evaluate LLMs on longitudinal clinical summarisation and
prediction using full patient trajectories from multi-modal EHRs. A key finding is
that longer context windows improve input \textit{integration} but do not reliably
improve clinical \textit{reasoning}. This connects to our capacity ceiling result:
the failure mode for small models on complex patients is not that the records do not
fit in context, but that the models fail to enumerate all active medications in their
output even when the information is present in the input. The reasoning bottleneck, not
the context window, is the limiting factor.

\subsection{Models Evaluated}

We evaluate Phi-3.5-mini \citep{abdin2024phi3}, Mistral-7B \citep{jiang2023mistral},
BioMistral-7B \citep{labrak2024biomistral}, and models from the Llama 3 family
\citep{llama3team2024}. Synthetic patient data is generated with Synthea
\citep{walonoski2018synthea}.
